\documentclass[11pt,a4paper]{article}
\usepackage[utf8]{inputenc}
\usepackage[T1]{fontenc}
\usepackage[french]{babel}
\usepackage{geometry}
\usepackage{array}
\usepackage{hyperref}
\usepackage{xcolor}
\usepackage{enumitem}
\usepackage{longtable}
\usepackage{titlesec}

\geometry{margin=2.5cm}
\hypersetup{
  colorlinks=true,
  linkcolor=blue,
  urlcolor=blue,
  pdftitle={Analyse technique de l'interface Smart Shepherd},
  pdfauthor={GitHub Copilot}
}

\title{Analyse technique de l'interface Smart Shepherd}
\author{Rapport technique}
\date{28 avril 2026}

\titleformat{\section}{\large\bfseries}{\thesection}{0.8em}{}
\titleformat{\subsection}{\normalsize\bfseries}{\thesubsection}{0.8em}{}

\begin{document}
\maketitle

\section{Vue d'ensemble}
L'application est un tableau de bord IoT pour la supervision de troupeaux. Elle combine une interface web moderne, un systeme de temps reel, une couche d'authentification, des visualisations de donnees, une carte interactive et un module d'IA.

\section{Technologies utilisees}
\subsection{Frontend}
\begin{itemize}[leftmargin=1.2cm]
  \item Framework principal : React 18.
  \item Bundler et environnement de dev : Vite.
  \item Navigation : React Router DOM.
  \item Etat global et cache serveur : Zustand et React Query.
  \item Requetes HTTP : Axios.
  \item Styling : Tailwind CSS, CSS custom properties, classes utilitaires et styles globaux.
  \item Cartographie : Leaflet, React Leaflet, Leaflet Draw, Leaflet Heat.
  \item Graphiques : Chart.js et React ChartJS 2.
  \item IoT et temps reel : MQTT cote client.
  \item PWA : vite-plugin-pwa.
  \item UI supplementaire : lucide-react, react-hot-toast.
\end{itemize}

\subsection{Backend}
\begin{itemize}[leftmargin=1.2cm]
  \item Runtime : Node.js.
  \item Framework API : Express.
  \item Base de donnees : MongoDB avec Mongoose.
  \item Authentification : JWT.
  \item Validation : Joi.
  \item Securite : Helmet, CORS, express-rate-limit.
  \item Temps reel : Socket.IO.
  \item IoT : MQTT.
  \item Journalisation : Winston.
  \item Chiffrement mot de passe : bcryptjs.
\end{itemize}

\section{Architecture du projet}
\subsection{Structure generale}
\begin{itemize}[leftmargin=1.2cm]
  \item \texttt{src/} : frontend React.
  \item \texttt{src/components/} : composants reutilisables, layout, widgets, IA et UI.
  \item \texttt{src/pages/} : pages metier du dashboard.
  \item \texttt{src/contexts/} : contexte d'authentification, theme et MQTT.
  \item \texttt{src/hooks/} : hooks personnalises et store Zustand.
  \item \texttt{src/services/} : clients API et services metier.
  \item \texttt{backend/} : API Express, routes, models, services et middleware.
  \item \texttt{ai-service/} : service IA Python separable.
  \item \texttt{firmware/} : code embarque ESP32.
\end{itemize}

\subsection{Evaluation de maintenabilite}
L'architecture est globalement scalable pour un prototype ou un produit en phase de croissance, car les couches sont deja segmentees. Cependant, il existe plusieurs fragilites :
\begin{itemize}[leftmargin=1.2cm]
  \item melange de fichiers \texttt{.jsx} et \texttt{.tsx}, ce qui complique la coherence typee ;
  \item dependances de presentation chargees de facon eager, ce qui rend le demarrage fragile ;
  \item store Zustand contenant des champs non parfaitement alignes avec son interface ;
  \item donnees backend parfois en memoire, donc non persistantes ;
  \item couplage fort entre certaines pages et leurs services.
\end{itemize}

\section{Description de l'interface}
\subsection{Ecran de connexion}
L'ecran de connexion propose un fond degrade, une carte translucide, des animations et un choix de compte de demonstration. L'approche est moderne et visuellement attractive, proche d'un style SaaS premium.

\subsection{Layout principal}
L'interface principale est composee de :
\begin{itemize}[leftmargin=1.2cm]
  \item une barre laterale avec navigation contextuelle ;
  \item un header avec nom de l'utilisateur, theme, notifications et statut live ;
  \item une zone de contenu centrale pour les pages metier ;
  \item un assistant IA flottant en bas a droite ;
  \item un systeme de toasts et d'alertes critiques.
\end{itemize}

\subsection{Pages fonctionnelles}
\begin{itemize}[leftmargin=1.2cm]
  \item Dashboard : KPIs, graphiques et alertes temps reel.
  \item Map : visualisation geographique des animaux et geofencing.
  \item Animals : liste et detail des animaux.
  \item Alerts : centre de notifications.
  \item Analytics : historique, tendances et export.
  \item AI Dashboard : predictions et anomalies IA.
  \item Hardware, Users, Settings : modules admin.
\end{itemize}

\section{Problemes detectes}
\subsection{Ecran blanc ou application bloquee}
Les causes les plus probables sont :
\begin{itemize}[leftmargin=1.2cm]
  \item import manquant pour \texttt{chartjs-plugin-annotation} dans \texttt{LabellingPage.jsx} ;
  \item import manquant pour \texttt{framer-motion} dans \texttt{AIInsightsFeed.tsx} ;
  \item crash serveur backend a cause de \texttt{@tensorflow/tfjs-node} absent ;
  \item verification d'authentification qui bloque le rendu tant que la session n'est pas resolue ;
  \item erreurs de chargement ou redirection dans les routes protegees.
\end{itemize}

\subsection{Probleme de contrat API}
Le frontend attend parfois un tableau brut, alors que le backend renvoie des objets enveloppes :
\begin{itemize}[leftmargin=1.2cm]
  \item \texttt{/api/sheep} renvoie \texttt{\{ sheep, pagination \}} ;
  \item \texttt{/api/notifications} renvoie \texttt{\{ notifications, unreadCount, pagination \}} ;
  \item le frontend de \texttt{AppLayout} verifie \texttt{Array.isArray(...)} et ignore donc les reponses valides.
\end{itemize}

\subsection{Incoherence des roles}
Le frontend utilise \texttt{SUPER\_ADMIN}, \texttt{ADMIN} et \texttt{FARMER}, alors que le backend valide \texttt{admin}, \texttt{operator} et \texttt{viewer}. Cette incoherence peut casser les menus, les droits et la logique d'autorisation.

\section{Analyse des performances}
\begin{itemize}[leftmargin=1.2cm]
  \item Le batching des mises a jour IoT dans Zustand est une bonne approche.
  \item Le chargement eager de nombreuses pages augmente le cout initial.
  \item Les imports lourds dans les routes peuvent degrader le temps de demarrage.
  \item Certains composants peuvent re-render plus souvent que necessaire si les selecteurs Zustand ne sont pas optimises.
  \item Les donnees en memoire cote backend ne sont pas adaptees a la production.
\end{itemize}

\section{Qualite UI/UX}
L'interface est deja dans une direction SaaS moderne : cartes, blur, gradients, badges live, navigation claire, analytics et cartes interactives. Les points a ameliorer sont :
\begin{itemize}[leftmargin=1.2cm]
  \item harmoniser la langue de l'interface ;
  \item rendre les etats vide, erreur et chargement plus explicites ;
  \item afficher les statuts backend, MQTT et synchronisation ;
  \item simplifier certaines vues denses ;
  \item charger les pages lourdes a la demande.
\end{itemize}

\section{Guide de debogage}
\begin{enumerate}[leftmargin=1.2cm]
  \item Ouvrir la console navigateur et verifier les erreurs d'import ou de rendu React.
  \item Lancer le build frontend pour detecter les modules manquants.
  \item Verifier le demarrage du backend et les logs Node.js.
  \item Tester \texttt{/api/health}, puis \texttt{/api/sheep} et \texttt{/api/notifications}.
  \item Confirmer que le token est present dans \texttt{localStorage}.
  \item Verifier que les roles backend et frontend correspondent.
  \item Inspecter le Network tab pour les erreurs 401, 404 et 500.
\end{enumerate}

\section{Plan d'amelioration}
\begin{itemize}[leftmargin=1.2cm]
  \item Ajouter les dependances manquantes ou supprimer les imports inutilises.
  \item Normaliser les reponses API et les types de donnees.
  \item Fractionner les grosses pages via \texttt{React.lazy}.
  \item Centraliser les types partagees frontend/backend.
  \item Remplacer les donnees in-memory par une persistence MongoDB.
  \item Unifier les roles utilisateur sur toute la stack.
  \item Renforcer les tests de build, d'API et de rendu.
\end{itemize}

\section{Conclusion}
Le projet a une base technique solide et une interface deja proche d'un produit SaaS moderne. Le principal point bloquant n'est pas le design, mais la fiabilite du pipeline : dependances manquantes, backend instable et contrat API incoherent. Une fois ces problemes corriges, le dashboard sera beaucoup plus stable, maintenable et facile a faire evoluer.

\end{document}
